\section*{\Huge\textcolor{tcd_blue}{Declaration}}

I declare that this dissertation has not been submitted as an exercise for a degree at this or any
other university and that it is entirely my own work.

I agree to deposit this dissertation in the University's open access institutional repository or
allow the Library to do so on my behalf, subject to Irish Copyright Legislation and Trinity College
Library conditions of use and acknowledgement.

I consent / do not consent to the examiner retaining a copy of the dissertation beyond the examining
period, should they so wish (EU GDPR May 2018).

\textbf{Use of generative artificial intelligence:} \textit{[Complete per TCD's current policy on
generative-AI use in research writing — state which tools were used, for which purposes (e.g.
literature-search assistance, drafting support, code assistance), and affirm that the research
design, system implementation decisions, and intellectual contributions are the author's own.]}

\vspace{3cm}

Signed:~\rule{5cm}{0.3pt}\hfill Date:~\rule{5cm}{0.3pt}

\newpage

\section*{\Huge\textcolor{tcd_blue}{Abstract}}

Extended-reality headsets are conventionally used to \textit{add} information to the world. This
dissertation investigates the opposite operation: using a consumer video-see-through (VST) headset to
\textit{subtract} by dimming, muting, or re-grading the visual periphery so that attention is guided towards,
and held on, to a user-chosen region of the real world. The work asks one question with two halves: can
subtractive attention guidance be delivered on consumer XR hardware, and, when it is delivered, does
it actually help human attention?

The technical half confronts a defining constraint of the Meta Quest 3 platform: the operating system
owns the passthrough layer, which applications may neither read nor spatially modify. The
dissertation contributes a three-tier design space of diminished reality (DR) under this compositing
constraint, and a working reference implementation spanning it: a world-locked \textit{focus window} rendered as an absence in an overlay sphere, through which native passthrough passes at full quality, with two peripheral treatments  carried to confirmatory evaluation, SignPop (detection-gated salience re-grading of the camera feed) and Hard Dark (full peripheral occlusion), plus motion-coupled suppression for
comfort. A structured prior-art search found no published or shipped precedent for the system's
hybrid composite of native passthrough and a shader-processed camera feed. A technical evaluation
plan (frame cost, latency, legibility across rendering paths, and the gap between live and oracle
object detection) is specified.

The human half contributes a pre-registered, fully within-subjects user study (N = 20 target) built
to be falsifiable rather than merely hopeful: a workstation distractor-suppression paradigm carrying
the single primary hypothesis, a driving-scene marker-detection paradigm with an equivalence-bounded
situational-awareness safety test, and a pre-committed decision table under which every outcome,
positive, negative, or qualified, is a conclusive result. Recruitment is ongoing, and eleven
participant pairs have been tested so far. Interim results from those eleven show a 3.65-point
accuracy gain on the workstation task, no loss of peripheral awareness, and a detection benefit
concentrated in the eccentricity band the system's own detection-gated design predicts. These are
reported and discussed as interim findings throughout, not as the final confirmatory result.

\newpage

\section*{\Huge\textcolor{tcd_blue}{Acknowledgements}}

\textit{[ACKNOWLEDGEMENTS — supervisor, pilot participants, School of Computer Science and Statistics,
family/friends, as appropriate.]}

\newpage

\section*{\Huge\textcolor{tcd_blue}{Summary of Contributions}}

This dissertation makes four contributions:

\begin{itemize}
\item \textbf{C1 — Design space.} A taxonomy of diminished-reality capability under compositing constraints on
  consumer video-see-through hardware (three access tiers: the OS-owned passthrough layer, overlay
  compositing, and camera re-render), with the achievable effects, ceilings, and costs of each tier
  characterised empirically (Chapter 3).

\item \textbf{C2 — Reference implementation.} An open, working multi-mode DR attention-guidance system on
  Quest 3 passthrough,  the world-locked focus-window-as-absence architecture, world-direction fusion
  sampling for the monocular camera feed, and motion-coupled comfort suppression, including
  documented failure modes and dead ends (Chapter 4).

\item \textbf{C3 — Technical evaluation.} A benchmark plan covering per-mode rendering cost, latency,
  legibility across rendering paths, and the measured distance between live on-device perception and
  the oracle-perception assumption used in the user study, together with preliminary results from the
  offline detection pipeline (Chapter 5).

\item \textbf{C4 — Pre-registered user study, with interim results.} A falsifiable, within-subjects
  evaluation design (N = 20 target) of peripheral diminishment under controlled distraction, with
  equivalence-bounded safety hypotheses and a pre-committed decision table covering positive,
  negative, and qualified outcomes, together with an interim analysis of the eleven participant pairs
  collected so far (Chapter 6).
\end{itemize}

\newpage
